% ============================================================================
\chapter{\'Ev\'enements, formulaires et interaction utilisateur}
\label{ch:js-events}
% ============================================================================

\begin{objectif}
\begin{itemize}
  \item Ma\^itriser les trois phases de propagation et l'art de d\'el\'eguer un \'ecouteur.
  \item Exploiter les options d'\texttt{addEventListener} (\texttt{signal}, \texttt{once},
        \texttt{capture}, \texttt{passive}).
  \item Valider un formulaire c\^ot\'e client \emph{et} comprendre pourquoi ce n'est qu'un confort.
  \item Envoyer des fichiers, g\'erer upload avec progression, et les \'ev\'enements clavier/souris.
  \item Construire les composants UI courants (menu, modale, autocomplete, toast).
\end{itemize}
\end{objectif}

\section{Le cycle de vie d'un \'ev\'enement}

Un clic sur un \verb|<span>| dans un \verb|<button>| dans une \verb|<form>| parcourt l'arbre
deux fois~: \textbf{capture} (du plus grand vers la cible), \textbf{target} (la cible),
\textbf{bulle} (de la cible vers le plus grand). La plupart du temps on \'ecoute en phase de
bulle, c'est pour \c{c}a que la d\'el\'egation fonctionne.

\begin{codejs}
form.addEventListener(
  "submit",
  (e) => {
    e.preventDefault();                  // empeche le comportement par defaut (rechargement)
    console.log(e.type, e.target, e.currentTarget);   // currentTarget = l'element quiecoute
    console.log(e.clientX, e.clientY, e.timeStamp);
  },
  { signal: abort.signal }               // la meilleure facon de nettoyer
);

const options = {
  capture: true,        // ecouter en phase descendante
  once: true,           // se desabonne apres le premier declenchement
  passive: true,        // promet de ne pas appeler preventDefault -> scroll plus fluide
  signal: new AbortController().signal,   // desabonnement en masse
};
window.addEventListener("wheel", onWheel, { passive: true });

// Stopper la propagation vs empecher l'action par defaut : DEUX choses differentes
e.stopPropagation();        // les parents ne seront pas notifies (mais l'action a lieu)
e.stopImmediatePropagation(); // ni les autres ecouteurs du MEME element
e.preventDefault();         // l'action par defaut est annulee (la propagation continue)

// Evenements sans bulle : focus, blur, mouseenter, load, scroll (sur element)
// => pour les deleguer, ecouter en capture :
conteneur.addEventListener("focusin", onFocus);      // focusin/focusout, eux, bullent
conteneur.addEventListener("focus", onFocus, true);  // capture : marche aussi
\end{codejs}

\section{La d\'el\'egation~: un \'ecouteur pour mille lignes}

C'est LE motif qui fait la diff\'erence entre un script qui vieillit bien et un script
qu'on r\'e\'ecrit chaque fois qu'on ajoute un \'el\'ement.

\begin{codejs}
// ANTI-PATTERN : attacher un ecouteur par ligne, a reattacher a chaque rendu
document.querySelectorAll(".suppr").forEach((b) => b.addEventListener("click", supprimer));

// DELEGATION : un seul ecouteur sur un conteneur stable
document.getElementById("liste").addEventListener("click", (e) => {
  const boutonSuppr = e.target.closest("[data-action=supprimer]");
  if (boutonSuppr) return supprimer(boutonSuppr.closest("li").dataset.id);

  const coche = e.target.closest("[data-action=cocher]");
  if (coche) return cocher(coche.closest("li").dataset.id);

  const lien = e.target.closest("a");
  if (lien) { e.preventDefault(); naviguer(lien.href); }
});
\end{codejs}

Un utilitaire de bus d'\'ev\'enements, \`a copier dans ton projet (il remplace avantageusement
les r\'ef\'erences crois\'ees entre modules)~:

\begin{codejs}
export function creerBus() {
  const ecouteurs = new Map();
  return {
    on(type, fn, options = {}) {
      if (!ecouteurs.has(type)) ecouteurs.set(type, new Set());
      ecouteurs.get(type).add(fn);
      return () => this.off(type, fn);
    },
    once(type, fn) { const off = this.on(type, (p) => { off(); fn(p); }); return off; },
    off(type, fn) { ecouteurs.get(type)?.delete(fn); },
    emit(type, detail) {
      for (const fn of ecouteurs.get(type) ?? []) {
        try { fn(detail); } catch (err) { console.error(`ecouteur "${type}"`, err); }
      }
      if (typeof CustomEvent === "function") {
        document.dispatchEvent(new CustomEvent(type, { detail, bubbles: true }));
      }
    },
    clear() { ecouteurs.clear(); },
  };
}
// Evenements personnalises, la norme DOM :
document.body.dispatchEvent(new CustomEvent("panier:maj", { detail: { total: 42 }, bubbles: true, composed: true }));
document.body.addEventListener("panier:maj", (e) => afficher(e.detail.total));
\end{codejs}

\section{Nettoyer : le pi\`ege n\${}^{\text{o}}\$1 des SPA}

Un \'ecouteur ajout\'e et jamais retir\'e = fuite de m\'emoire + rappels fant\^omes sur des \'el\'ements
d\'etruits. Trois strat\'egies, de la plus simple \`a la plus propre.

\begin{codejs}
// 1. once : pour un evenement unique
btn.addEventListener("click", demarrer, { once: true });

// 2. AbortController : le standard moderne, et le plus pratique
const ctrl = new AbortController();
window.addEventListener("resize", surResize, { signal: ctrl.signal });
document.addEventListener("keydown", surTouche, { signal: ctrl.signal });
form.addEventListener("submit", surEnvoi, { signal: ctrl.signal, passive: false });
function detruire() { ctrl.abort(); }       // retire TOUT, en une ligne

// 3. Un mini-cycle de vie pour tes propres composants
function creerComposant(racine, { montar, demonter }) {
  const ctrl = new AbortController();
  montar(racine, ctrl.signal);
  return {
    destroy() { ctrl.abort(); demonter?.(racine); racine.replaceChildren(); },
  };
}
\end{codejs}

\section{Formulaires---lecture, \'ecriture, validation}

\begin{codejs}
const form = document.querySelector("form");

// Lecture moderne : FormData
const data = new FormData(form);
const obj = Object.fromEntries(data.entries());           // { nom: "...", qte: "2" }
const roles = data.getAll("role");                         // checkbox multiple -> tableau
const fichier = data.get("avatar");                        // File | null
data.set("nom", "Nour"); data.append("tag", "admin");
// Renvoyer un objet typ\'e, propre a envoyer en JSON :
const lire = (form) => {
  const o = Object.fromEntries(new FormData(form).entries());
  for (const [k, v] of Object.entries(o)) {
    if (v === "on") o[k] = true;
    else if (v === "off") o[k] = false;
    else if (v !== "" && !Number.isNaN(Number(v)) && form.elements[k]?.type === "number") o[k] = Number(v);
  }
  return o;
};

// Validation native (le navigateur la fait deja, et gratuitement)
const champ = form.elements.email;
champ.checkValidity();          // booleen
champ.validity.tooLong; champ.validationMessage;
form.reportValidity();          // affiche les bulles, met le focus sur le 1er invalide
champ.setCustomValidity("Choisis un email @exemple.com");   // remplace le message
champ.setCustomValidity("");     // ne pas oublier de le reinitialiser !

// Les attributs HTML qui font 80 % du travail :
// required, type, pattern, minlength/maxlength, min/max/step, inputmode, autocomplete,
// autocapitalize, multiple, accept, list (datalist)
\end{codejs}

\subsection{Validation sur mesure, r\'eutilisable}

\begin{codejs}
const regles = {
  requis: (v) => (v === null || v === undefined || String(v).trim() === "") && "Champ obligatoire",
  email: (v) => v && !/^[^@\s]+@[^@\s]+\.[a-z]{2,}$/i.test(v) && "Adresse email invalide",
  min: (n) => (v) => [...String(v ?? "")].length < n && `Au moins ${n} caracteres`,
  nombre: (min, max) => (v) => {
    const x = Number(v);
    if (v === "" || !Number.isFinite(x)) return "Nombre attendu";
    if (min !== undefined && x < min) return `Minimum ${min}`;
    if (max !== undefined && x > max) return `Maximum ${max}`;
  },
  mdpFort: (v) => {
    if (!v) return;
    const manques = [];
    if (v.length < 12) manques.push("12 caracteres minimum");
    if (!/[A-Z]/.test(v)) manques.push("une majuscule");
    if (!/[a-z]/.test(v)) manques.push("une minuscule");
    if (!/\d/.test(v)) manques.push("un chiffre");
    if (!/[^\w\s]/.test(v)) manques.push("un caractere special");
    return manques.length && `Mot de passe faible : ${manques.join(", ")}`;
  },
};

export function valider(form, schema) {
  const erreurs = {};
  for (const [nom, liste] of Object.entries(schema)) {
    const champ = form.elements[nom];
    const valeur = champ?.type === "checkbox" ? champ.checked : champ?.value;
    for (const regle of liste) {
      const message = typeof regle === "function" ? regle(valeur, form) : regles[regle]?.(valeur, form);
      if (message) { erreurs[nom] = message; break; }
    }
  }
  return { ok: Object.keys(erreurs).length === 0, erreurs };
}

// Branchement : invalider a la sortie, revalider a la frappe une fois l'erreur affichee
const champ = form.elements.email;
const afficherErreur = (msg) => {
  const zone = champ.closest(".champ");
  zone.classList.toggle("invalide", Boolean(msg));
  zone.querySelector(".erreur").textContent = msg ?? "";
  champ.setAttribute("aria-invalid", String(Boolean(msg)));
  champ.setAttribute("aria-describedby", msg ? "err-email" : "");
};
champ.addEventListener("blur", () => afficherErreur(regles.email(champ.value)));
champ.addEventListener("input", () => champ.closest(".champ").classList.contains("invalide")
  && afficherErreur(regles.email(champ.value)));

form.addEventListener("submit", async (e) => {
  e.preventDefault();
  const { ok, erreurs } = valider(form, {
    email: ["requis", "email"],
    age: [regles.nombre(18, 120)],
    mdp: [regles.mdpFort],
  });
  if (!ok) {
    const premier = form.elements[Object.keys(erreurs)[0]];
    afficherErreur(erreurs.email);
    premier?.focus();
    return;
  }
  await soumettre(Object.fromEntries(new FormData(form)));
});
\end{codejs}

\begin{piege}
Un formulaire doit \^etre soumis par un vrai \verb|<button type="submit">| et r\'epondre \`a
\verb|Enter|~: c'est ce que la touche Entr\'ee, les lecteurs d'\'ecran et les gestionnaires de mots
de passe attendent. Un simple \verb|onclick| sur une \verb|<div>| casse tout cela~;
la validation native ne se d\'eclenche m\^eme pas.
\end{piege}

\section{Fichiers : s\'election, lecture, envoi, progression}

\begin{codejs}
const input = document.querySelector('input[type=file]');

input.addEventListener("change", async () => {
  for (const file of input.files) {
    // Garde-fous cote client (le serveur refait TOUJOURS les siens)
    if (file.size > 5 * 1024 ** 2) { alert(`${file.name} depasse 5 Mo`); continue; }
    if (!/^image\/(png|jpe?g|webp|avif)$/.test(file.type)) continue;

    const url = URL.createObjectURL(file);            // lecture immediate pour preview
    document.querySelector("#apercu").src = url;

    const texte = await file.text();                  // File -> string (Promise)
    const buffer = await file.arrayBuffer();          // binaire (hash, analyse)
    const hash = [...new Uint8Array(await crypto.subtle.digest("SHA-256", buffer))]
      .map((b) => b.toString(16).padStart(2, "0")).join("");

    const dataUrl = await new Promise((res, rej) => {
      const r = new FileReader();
      r.onload = () => res(r.result);
      r.onerror = () => rej(r.error);
      r.readAsDataURL(file);
    });
    URL.revokeObjectURL(url);                         // liberer la memoire !
    console.log(file.name, file.size, file.type, file.lastModified, hash.slice(0, 12));
  }
});

// Envoi : multipart via FormData, le navigateur met le Content-Type tout seul
async function envoyer(form) {
  const fd = new FormData(form);
  const ctrl = new AbortController();
  const xhr = new XMLHttpRequest();           // seul moyen simple d'avoir la progression
  xhr.open("POST", "/upload");
  xhr.upload.onprogress = (e) => {
    if (e.lengthComputable) barre.style.width = `${(e.total * 100 / e.loaded).toFixed(0)}`;
  };
  xhr.send(fd);
}

// Upload en chunk (gros fichiers) : decoupage, envoi parallele, reprise
async function uploaderParMorceaux(file, taille = 5 * 1024 ** 2) {
  const sessions = await fetch("/upload/init", {
    method: "POST", headers: { "content-type": "application/json" },
    body: JSON.stringify({ nom: file.name, taille: file.size, morceaux: Math.ceil(file.size / taille) }),
  }).then((r) => r.json());
  const envois = [];
  for (let i = 0; i * taille < file.size; i++) {
    const morceau = file.slice(i * taille, (i + 1) * taille);
    envois.push(fetch(`/upload/${sessions.id}/${i}`, { method: "PUT", body: morceau }));
    if (envois.length >= 3) await Promise.race(envois);       // limiter la concurrence
  }
  await Promise.all(envois);
  return fetch(`/upload/${sessions.id}/complete`, { method: "POST" });
}

// Telecharger un fichier genere cote serveur, sans quitter la page
const reponse = await fetch("/export.csv");
const blob = await reponse.blob();
const a = Object.assign(document.createElement("a"), {
  href: URL.createObjectURL(blob), download: "export.csv",
});
a.click(); URL.revokeObjectURL(a.href);
\end{codejs}

\section{Clavier, souris, glisser, tactile}

\begin{codejs}
// Clavier : toujours tester la touche + les modificateurs, et menager les accessifs
window.addEventListener("keydown", (e) => {
  if (e.target.matches("input, textarea, [contenteditable]")) {
    if (!(e.ctrlKey || e.metaKey)) return;      // ne pas capturer les touches de saisie
  }
  const mod = e.metaKey || e.ctrlKey ? "Mod+" : "";
  const combo = `${mod}${e.shiftKey ? "Shift+" : ""}${e.key.toLowerCase()}`;
  const raccourcis = {
    "Mod+k": () => ouvrirPalette(),
    "Mod+Enter": () => soumettre(),
    "Mod+Shift+p": () => basculerMode(),
    Escape: () => fermerTout(),
    "?": () => ouvrirAide(),
  };
  const fn = raccourcis[combo];
  if (fn) { e.preventDefault(); fn(); }
});

// Souris : coalescer le mouvement (60+ evenements par frame sinon)
let last = { x: 0, y: 0 };
el.addEventListener("pointermove", (e) => {
  const [evt] = e.getCoalescedEvents?.() ?? [e];     // regroupe les evements perdus
  last = { x: evt.clientX, y: evt.clientY };
}, { passive: true });

// Pointer events : UN seul modele pour souris + stylet + doigt (oublie touch*)
const points = new Map();
el.addEventListener("pointerdown", (e) => {
  el.setPointerCapture(e.pointerId);
  points.set(e.pointerId, { x: e.clientX, y: e.clientY });
});
el.addEventListener("pointermove", (e) => points.get(e.pointerId) && pince(e, points));
el.addEventListener("pointerup", (e) => points.delete(e.pointerId));

// Double-clic, clic droit, molette, drag natif
el.addEventListener("dblclick", zoom);
el.addEventListener("contextmenu", (e) => { e.preventDefault(); menu.positionner(e.clientX, e.clientY); });
el.addEventListener("wheel", (e) => { if (e.ctrlKey) { e.preventDefault(); zoomPar(e.deltaY); } }, { passive: false });
el.draggable = true;
el.addEventListener("dragstart", (e) => e.dataTransfer.setData("text/plain", el.dataset.id));
zone.addEventListener("dragover", (e) => { e.preventDefault(); zone.classList.add("sur"); });
zone.addEventListener("drop", (e) => {
  e.preventDefault();
  deplacer(e.dataTransfer.getData("text/plain"), indexApres(e));
});
// Drag de fichiers depuis le bureau
window.addEventListener("dragover", (e) => e.preventDefault());
window.addEventListener("drop", (e) => { e.preventDefault(); traiter(e.dataTransfer.files); });
\end{codejs}

\section{Quatre composants UI, sans biblioth\`eque}

\begin{codejs}
// 1. Menu deroulant accessible (fleches, Echap, clic exterieur, focus)
function menu(bouton, liste) {
  bouton.addEventListener("click", () => {
    const ouvert = liste.hasAttribute("hidden");
    liste.toggleAttribute("hidden", !ouvert);
    bouton.setAttribute("aria-expanded", String(ouvert));
    if (ouvert) liste.querySelector('[role="menuitem"]')?.focus();
  });
  liste.addEventListener("keydown", (e) => {
    const items = [...liste.querySelectorAll('[role="menuitem"]:not([disabled])')];
    const i = items.indexOf(document.activeElement);
    const bouger = (n) => { e.preventDefault(); items[(i + n + items.length) % items.length].focus(); };
    if (e.key === "ArrowDown") bouger(1);
    else if (e.key === "ArrowUp") bouger(-1);
    else if (e.key === "Home") bouger(-i);
    else if (e.key === "End") bouger(items.length - 1 - i);
    else if (e.key === "Escape") { bouton.click(); bouton.focus(); }
  });
  document.addEventListener("click", (e) => {
    if (!e.target.closest("[data-menu]") && !liste.contains(e.target)) {
      liste.setAttribute("hidden", ""); bouton.setAttribute("aria-expanded", "false");
    }
  });
}

// 2. Toast (notification auto-detruite)
const toasts = (conteneur => ({
  show(message, { type = "info", duree = 4000 } = {}) {
    const t = document.createElement("div");
    t.className = `toast toast--${type}`; t.setAttribute("role", "status"); t.textContent = message;
    conteneur.append(t);
    t.addEventListener("click", () => t.remove());
    setTimeout(() => { t.classList.add("out"); t.addEventListener("transitionend", () => t.remove(), { once: true }); }, duree);
  },
}))(document.getElementById("toasts"));
toasts.show("Enregistre", { type: "succes" });

// 3. Autocompletion (combobox) avec debounce + annulation + navigation clavier
function autocomplete(input, liste, chercher) {
  let ctrl = null, index = -1;
  const fermer = () => { liste.setAttribute("hidden", ""); index = -1; };
  input.addEventListener("input", debounce(async () => {
    ctrl?.abort(); ctrl = new AbortController();
    const q = input.value.trim();
    if (q.length < 2) return fermer();
    let items;
    try { items = await chercher(q, ctrl.signal); } catch (e) { if (e.name !== "AbortError") throw e; return; }
    if (!items.length) return fermer();
    liste.replaceChildren(...items.map((it, i) => {
      const li = document.createElement("li");
      li.setAttribute("role", "option"); li.textContent = it.label; li.dataset.value = it.value;
      li.id = `ac-${i}`;
      return li;
    }));
    liste.removeAttribute("hidden");
    input.setAttribute("aria-expanded", "true");
  }, 220));
  input.addEventListener("keydown", (e) => {
    const opts = [...liste.querySelectorAll("[role=option]")];
    if (e.key === "ArrowDown" || e.key === "ArrowUp") {
      e.preventDefault();
      index = (index + (e.key === "ArrowDown" ? 1 : -1) + opts.length) % opts.length;
      opts.forEach((o, i) => o.setAttribute("aria-selected", String(i === index)));
      input.setAttribute("aria-activedescendant", opts[index]?.id ?? "");
    } else if (e.key === "Enter" && index >= 0) {
      e.preventDefault(); input.value = opts[index].dataset.value; fermer();
    } else if (e.key === "Escape") fermer();
  });
  liste.addEventListener("click", (e) => {
    const o = e.target.closest("[role=option]"); if (!o) return;
    input.value = o.dataset.value; fermer(); input.focus();
  });
}

// 4. Scroll infini (la seule bonne impl, sans listener de scroll)
function infini(sentinelle, charger) {
  let enCours = false;
  const io = new IntersectionObserver(async ([e]) => {
    if (!e.isIntersecting || enCours) return;
    enCours = true;
    const fini = await charger();
    enCours = false;
    if (fini) io.disconnect();
  }, { rootMargin: "400px" });
  io.observe(sentinelle);
  return () => io.disconnect();
}
\end{codejs}

\begin{exercice}
\begin{enumerate}
  \item Reprends la liste de courses et remplace tous les \'ecouteurs par un seul syst\`eme de
        d\'el\'egation + un \verb|AbortController| global.
  \item \'Ecris un formulaire d'inscription complet~: nom, email, mot de passe (jauge de force
        en temps r\'eel), confirmation, conditions (case \`a cocher). Validation \`a la sortie du champ,
        messages sous le champ, \verb|aria-invalid|, focus sur la premi\`ere erreur.
  \item Construis un s\'electeur de date \og ~calendrier~ \fg{}~: grille 7\$\times\$6, navigation mois,
        \keys{\$\leftarrow\$} \keys{\$\rightarrow\$}, plage interdite, valeur au format \verb|YYYY-MM-DD| dans un champ cach\'e.
  \item Ajoute un glisser-d\'eposer de fichiers sur toute la page, avec aper\c{c}u des images,
        rejet des types interdits et affichage de la taille.
  \item \'Ecris une barre de progression de t\'el\'echargement r\'ealiste~: \verb|fetch| +
        \verb|ReadableStream| pour compter les octets re\c{c}us (sans \verb|XMLHttpRequest|).
  \item Fais un raccourci clavier global \og ~palette de commandes~ \fg{} (Ctrl+K) avec recherche
        floue, haut/bas, Entr\'ee, et focus pi\'eg\'e.
\end{enumerate}
\end{exercice}

\begin{ressources}
\begin{itemize}
  \item MDN --- \'Ev\'enements --- \url{https://developer.mozilla.org/fr/docs/Web/API/EventTarget/addEventListener}
  \item MDN --- Constraint Validation API --- \url{https://developer.mozilla.org/fr/docs/Web/API/HTMLConstraintValidation}
  \item \emph{Form Validation Design Patterns} (GOV.UK) --- \url{https://design-system.service.gov.uk/patterns/validation/}
  \item \link{a11y-collective / ARIA Authoring Practices}{https://www.w3.org/WAI/ARIA/apg/patterns/} (mod\`eles de composants accessibles)
  \item \link{File Picking, les 10 commandes}{https://web.dev/file-system-access/}
\end{itemize}
\end{ressources}

\index{evenements}\index{formulaire}\index{validation}\index{accessibilite}
